\documentclass[a4paper]{article}
\usepackage{amsmath, amsthm, amsfonts, amssymb, mathrsfs}
\usepackage[all]{xy}
\usepackage{color}
\usepackage{authblk}
\usepackage{afterpage}
\usepackage{xcolor}
\usepackage[backref=page]{hyperref}
\usepackage{aliascnt}
\usepackage[capitalise]{cleveref}
\usepackage[a4paper, bindingoffset=0.2in, left=1.3in, right=1.3in, top=1.4in, bottom=1.5in, footskip=.5in]{geometry}
\usepackage{leftindex}
\usepackage{extarrows}
\usepackage{float}
\usepackage{graphicx}
\usepackage{tikz}
\usetikzlibrary{patterns}

\theoremstyle{plain}
\newtheorem{theorem}{Theorem}[section]

\newaliascnt{proposition}{theorem}
\newtheorem{proposition}[proposition]{Proposition}
\aliascntresetthe{proposition}
\crefname{proposition}{Proposition}{Propositions}

\newaliascnt{corollary}{theorem}
\newtheorem{corollary}[corollary]{Corollary}
\aliascntresetthe{corollary}
\crefname{corollary}{Corollary}{Corollaries}

\newaliascnt{lemma}{theorem}
\newtheorem{lemma}[lemma]{Lemma}
\aliascntresetthe{lemma}
\crefname{lemma}{Lemma}{Lemmas}

\theoremstyle{definition}
\newaliascnt{definition}{theorem}
\newtheorem{definition}[definition]{Definition}
\aliascntresetthe{definition}
\crefname{definition}{Definition}{Definitions}

\newaliascnt{example}{theorem}
\newtheorem{example}[example]{Example}
\aliascntresetthe{example}
\crefname{example}{Example}{Examples}

\newtheorem{remark}{Remark}

\newaliascnt{notation}{theorem}
\newtheorem{notation}[notation]{Notation}
\aliascntresetthe{notation}
\crefname{notation}{Notation}{Notations}

\theoremstyle{plain}
\newtheorem*{theorem*}{Theorem}
\newtheorem*{proposition*}{Proposition}
\newtheorem*{corollary*}{Corollary}
\newtheorem*{lemma*}{Lemma}
\theoremstyle{definition}
\newtheorem*{definition*}{Definition}
\newtheorem*{remark*}{Remark}
\newtheorem*{notation*}{Notation}

\newcommand{\CC}{\mathbb{C}}
\newcommand{\PP}{\mathbb{P}}
\newcommand{\ZZ}{\mathbb{Z}}

\newcommand{\UUb}{\mathbb{U}}
\newcommand{\NN}{\mathbb{N}}
\newcommand{\OO}{\mathcal{O}}
\newcommand{\UU}{\mathcal{U}}
\newcommand{\DD}{\mathcal{D}}
\newcommand{\FF}{\mathcal{F}}
\newcommand{\EE}{\mathcal{E}}
\newcommand{\LL}{\mathcal{L}}
\newcommand{\RR}{\mathbb{R}}
\newcommand{\BB}{\mathcal{B}}
\newcommand{\KK}{\mathcal{K}}
\newcommand{\IIF}{\mathcal{I_F}}
\newcommand{\HH}{\mathcal{H}}
\newcommand{\C}{\mathcal{C}}
\newcommand{\II}{\mathscr{I}}
\newcommand{\JJ}{\mathscr{J}}
\newcommand{\pp}{\mathfrak{p}}

\DeclareMathAlphabet{\mathpzc}{OT1}{pzc}{m}{it}

\renewcommand{\ker}{\mathrm{Ker}}
\renewcommand{\ln}{\mathrm{ln}}
\newcommand{\id}{\mathrm{id}}
\newcommand{\hol}{\mathrm{Hol}}
\newcommand{\aut}{\mathrm{Aut}}
\newcommand{\im}{\mathrm{Im}}
\newcommand{\codim}{\mathrm{codim}}
\newcommand{\ann}{\mathrm{ann}}
\newcommand{\proj}{\mathrm{Proj}}
\newcommand{\ext}{\mathrm{Ext}}
\renewcommand{\hom}{\mathrm{Hom}}
\newcommand{\sing}{\mathrm{Sing}}
\newcommand{\reg}{\mathrm{Reg}}
\newcommand{\per}{\mathpzc{Per}(\omega)}
\newcommand{\kup}{\mathpzc{Kup}(\omega)}
\newcommand{\kupv}{\mathpzc{Kup}(\varpi)}
\newcommand{\wt}{\widetilde}
\newcommand{\defom}[1][]{
  \if\relax\detokenize{#1}\relax
    \omega_\varepsilon
  \else
    \leftindex_{#1}\omega_\varepsilon
  \fi
}
\usepackage{xparse}

\NewDocumentCommand{\unfom}{oo}{%
  \IfNoValueTF{#1}{%
    \IfNoValueTF{#2}{%
      \widetilde{\omega}_\varepsilon                            % no indices
    }{%
      \leftindex^{#2}{\widetilde{\omega}}_\varepsilon          % top-left only
    }%
  }{%
    \IfNoValueTF{#2}{%
      \leftindex_{#1}{\widetilde{\omega}}_\varepsilon          % bottom-left only
    }{%
      \leftindex^{#2}_{#1}{\widetilde{\omega}}_\varepsilon     % both
    }%
  }%
}

% \title{Holonomy is stable under unfoldings}
% \author[1]{Ariel Molinuevo\thanks{The author was fully supported by Universidade Federal do Rio de Janeiro, Brasil.}}
% \author[2]{Federico Quallbrunn\thanks{The author was fully supported by CONICET, Argentina.}}
% \affil[1]{\small Instituto de Matem\'atica, Universidade Federal do Rio de Janeiro, Brazil.}
% \affil[2]{\small Departamento de Matem\'atica, Facultad de Ciencias Exactas y Naturales, Universidad de Buenos Aires, Argentina.}
% \date{}

\newcommand\blankpage{
    \null
    \thispagestyle{empty}
    \addtocounter{page}{-1}
    \newpage}


% \author{}

\begin{document}

\rotatebox{-90}{%
\begin{tikzpicture}[scale=1.2]
  % Parameters
  \def\a{2}   % half-width at top
  \def\b{0.6} % perspective foreshortening (y-axis)
  \def\h{3}   % height of paraboloid

  % Floor plane rectangle at z=0 (appears on the left after rotation)
  \fill[gray!12] ({-\a-0.5}, {-\b*\a-0.4}) rectangle ({\a+0.5}, {\b*\a+0.4});
  \draw[gray!60, thin] ({-\a-0.5}, {-\b*\a-0.4}) rectangle ({\a+0.5}, {\b*\a+0.4});

  % Projection of paraboloid onto the floor
  \fill[blue!12] (0,0) ellipse[x radius=\a, y radius={\b*\a}];
  \pgfmathsetmacro{\rsmall}{\a*sqrt(1/\h)}
  \draw[thick, black!70!black] ({cos(40)*\a}, {\b*\a*sin(40)}) arc[x radius=\a, y radius={\b*\a}, start angle=40, end angle=-220];
  \draw[gray!50, thick, dashed] ({cos(40)*\a}, {\b*\a*sin(40)}) arc[x radius=\a, y radius={\b*\a}, start angle=40, end angle=140];

  % Left and right bounding parabolas
  \draw[thick, blue!70!black]
      plot[domain=0:\h*85/100, samples=40]
          ({-\a*sqrt(\x/\h)}, \x);
  \draw[thick, blue!70!black, dashed]
      plot[domain=\h*85/100:\h, samples=20]
          ({-\a*sqrt(\x/\h)}, \x);
  \draw[thick, blue!70!black]
      plot[domain=0:\h*85/100, samples=40]
          ({\a*sqrt(\x/\h)}, \x);
  \draw[thick, blue!70!black, dashed]
      plot[domain=\h*85/100:\h, samples=20]
          ({\a*sqrt(\x/\h)}, \x);

  % Top ellipse (solid front, dashed back)
  \draw[thick, blue!70!black]
      (\a, \h) arc[x radius=\a, y radius=\b*\a, start angle=0, end angle=-180];
  \draw[thick, blue!70!black]
      (\a, \h) arc[x radius=\a, y radius=\b*\a, start angle=0, end angle=180];

  % Front and back meridians of first paraboloid
  
  
  
  
  
  
  
  
  
  
  
  
  
  
  
  
  
  
  
  
  

  % Above floor (z=0 to z=h-1=2): solid up to 85%, dashed top 15%
  \draw[thick, blue!70!black]
      plot[domain=0:(\h-1)*75/100, samples=40]
          ({-\a*sqrt((\x+1)/\h)}, \x);
  \draw[thick, blue!70!black, dashed]
      plot[domain=(\h-1)*75/100:\h-1, samples=20]
          ({-\a*sqrt((\x+1)/\h)}, \x);
  \draw[thick, blue!70!black]
      plot[domain=0:(\h-1)*75/100, samples=40]
          ({\a*sqrt((\x+1)/\h)}, \x);
  \draw[thick, blue!70!black, dashed]
      plot[domain=(\h-1)*75/100:\h-1, samples=20]
          ({\a*sqrt((\x+1)/\h)}, \x);
  % Below floor (z=-1 to z=0): dashed
  \draw[thick, blue!70!black, dashed]
      plot[domain=-1:0, samples=30]
          ({-\a*sqrt((\x+1)/\h)}, \x);
  \draw[thick, blue!70!black, dashed]
      plot[domain=-1:0, samples=30]
          ({\a*sqrt((\x+1)/\h)}, \x);

  % Smaller inner circle on floor (drawn last to stay on top)
  \draw[thick, black] (\rsmall, 0) arc[x radius=\rsmall, y radius={\b*\rsmall}, start angle=0, end angle=-180];
  \draw[gray!50, thick, dashed] (\rsmall, 0) arc[x radius=\rsmall, y radius={\b*\rsmall}, start angle=0, end angle=180];

  % First paraboloid meridians (drawn after inner circle to appear on top)
  \draw[thick, blue!70!black]
      plot[domain=0:\h*55/100, samples=40]
          ({0}, {\x - \b*\a*sqrt(\x/\h)});
  \draw[thin, blue!70!black]
      plot[domain=\h*55/100:\h, samples=30]
          ({0}, {\x - \b*\a*sqrt(\x/\h)});
  \draw[thin, blue!50!black, dashed]
      plot[domain=\h*13/100:\h, samples=70]
          ({0}, {\x + \b*\a*sqrt(\x/\h)});

  % Front and back meridians of second paraboloid
  \draw[thick, blue!70!black]
      plot[domain=0:(\h-1)*75/100, samples=40]
          ({0}, {\x - \b*\a*sqrt((\x+1)/\h)});
  \draw[thick, blue!70!black, dashed]
      plot[domain=(\h-1)*75/100:\h-1, samples=20]
          ({0}, {\x - \b*\a*sqrt((\x+1)/\h)});
  \draw[thick, blue!70!black, dashed]
      plot[domain=-0.5:0, samples=30]
          ({0}, {\x - \b*\a*sqrt((\x+1)/\h)});
  \draw[thin, blue!50!black, dashed]
      plot[domain=(\h-1)*10/100:\h-1, samples=60]
          ({0}, {\x + \b*\a*sqrt((\x+1)/\h)});
  % Top ellipse of second paraboloid at z=h-1=2
  \draw[thick, blue!70!black]
      ({cos(25)*\a}, {\h-1+\b*\a*sin(25)}) arc[x radius=\a, y radius=\b*\a, start angle=25, end angle=-205];
  \draw[thick, blue!70!black, dashed]
      ({cos(25)*\a}, {\h-1+\b*\a*sin(25)}) arc[x radius=\a, y radius=\b*\a, start angle=25, end angle=180];

  % Vertex dot (drawn last to stay on top)
  \fill (0,0) circle (1.5pt);

\end{tikzpicture}%
}
\end{document}
